\documentclass[11pt,a4paper]{article}
% Shared preamble for RuPhO English problem statements (match T1 2026 style).
% Use: \input{latex/rupho_en_preamble.tex} after \documentclass.
\usepackage[margin=1in]{geometry}
\usepackage{amsmath,amssymb}
\usepackage{graphicx}
\usepackage{float}
\usepackage{enumitem}
\usepackage{microtype}
\usepackage{caption}
\usepackage{tabularx}
\usepackage{hyperref}

\captionsetup{font=small,labelfont=bf,skip=6pt}

\setlist[itemize]{leftmargin=1.2cm}
\setlist[enumerate]{leftmargin=1.2cm}

% One outer box per subproblem: [id | statement | points] (no inner rules)
\newcommand{\problembox}[3]{%
  \par\addvspace{0.42em}%
  \noindent
  \setlength{\fboxsep}{7.5pt}%
  \setlength{\fboxrule}{0.95pt}%
  \fbox{%
    \begin{minipage}{\dimexpr\linewidth-2\fboxsep-2\fboxrule\relax}
    \begin{tabularx}{\linewidth}{@{}>{\raggedright\arraybackslash\bfseries}p{2.1em} @{\hspace{0.9em}} >{\raggedright\arraybackslash}X @{\hspace{0.9em}} >{\raggedleft\arraybackslash\bfseries}p{2.4em}@{}}
    #1 & #2 & #3
    \end{tabularx}
    \end{minipage}%
  }%
  \par\addvspace{0.32em}%
}


\title{\textbf{Curling}}
\author{\normalsize Y21 (2021) --- English translation}
\date{}
\begin{document}
\maketitle

A thin ring of mass $m$ and radius $r$ is placed on a table with friction coefficient $\mu$ and gravitational acceleration $g$. The ring was given a linear velocity $v$ (in the direction $x$) and an angular velocity $\omega$ (counterclockwise). The ring moves in a straight line, rotating around its own axis, and eventually stops.  In the future we will try to find the quantities $v(t)$, $\omega(t)$ and some symmetry between them. In all subsequent paragraphs, it is assumed that the action of the support reaction on the ring is uniformly distributed.

You may need the following formula: $$ \textbackslash\{\}sqrt\{\textbackslash\{\}frac\{1 - \textbackslash\{\}sin \textbackslash\{\}theta\}\{2\}\} = \textbackslash\{\}Bigg| \textbackslash\{\}sin \textbackslash\{\}left( \{\textbackslash\{\}frac\{\textbackslash\{\}theta\}\{2\}\} - \textbackslash\{\}frac\{\textbackslash\{\}pi\}\{4\} \textbackslash\{\}right) \textbackslash\{\}Bigg| $$

\section*{Equations of motion}

\problembox{A1}{Show that the total force acting on the ring is given by: $$ \textbackslash\{\}vec F_\{tot\} = - \textbackslash\{\}mu mg \textbackslash\{\}cdot f \textbackslash\{\}left( \textbackslash\{\}frac\{v(t)\}\{\textbackslash\{\}omega (t)\textbackslash\{\} r\} \textbackslash\{\}right) \textbackslash\{\}hat x, $$ где $$ f(a) = \textbackslash\{\}frac\{1\}\{2 \textbackslash\{\}pi\} \textbackslash\{\}int \textbackslash\{\}limits_0^\{2 \textbackslash\{\}pi\} \textbackslash\{\}frac\{a - \textbackslash\{\}sin \textbackslash\{\}theta\}\{\textbackslash\{\}sqrt\{ 1 + a^2 - 2 a \textbackslash\{\}sin \textbackslash\{\}theta\}\} d \textbackslash\{\}theta $$}{0.50}

\problembox{A2}{Show that the total moment acting on the ring is: $$ \textbackslash\{\}tau_\{tot\} = - \textbackslash\{\}mu mg r f \textbackslash\{\}left(\textbackslash\{\}frac\{\textbackslash\{\}omega(t) r\}\{v(t)\} \textbackslash\{\}right) $$}{0.50}

\problembox{A3}{Prove that the equations of motion have the form: $$ \textbackslash\{\}dot v = - \textbackslash\{\}mu g\textbackslash\{\}cdot  f \textbackslash\{\}left( \textbackslash\{\}frac\{v\}\{\textbackslash\{\}omega r\} \textbackslash\{\}right) \textbackslash\{\}\textbackslash\{\} \textbackslash\{\}dot \textbackslash\{\}omega r = - \textbackslash\{\}mu g \textbackslash\{\}cdot f \textbackslash\{\}left( \textbackslash\{\}frac\{ \textbackslash\{\}omega r\}\{v\} \textbackslash\{\}right) $$}{0.10}

\section*{Primary Research}

We have obtained a system of differential equations relating $v(t)$ and $\omega (t)$. In this part we want to find interesting physical aspects of this situation, for which we need to make non-standard mathematical calculations. Let's start by considering the properties of the function $f(a)$:

\problembox{B1}{Prove: a) $ f(0) = 0, \: f(1) = \dfrac{2}{\pi}, \: f(\infty) = 1$ b) $ f(a) $ strictly increases as $a \geqslant 0 $}{0.50}

\problembox{B2}{Let us consider the behavior of the parameter $a(t)  = \dfrac{v(t)}{\omega (t) r} $. Show what happens to $a(t)$ (increases/decreases/remains unchanged) in each of the following cases: a) at some moment $a(t) = 1$ b) at some moment $a(t) < 1$ c) at some moment $a(t) > 1$}{0.30}

\problembox{B3}{Draw qualitatively on a graph, the axes of which are $v$ and $\omega r$, trajectories that display different movements of the ring, that is, given $v_0$ and $\omega_0 r$, draw how they will change over time. It is necessary to draw at least one trajectory for each point of the previous task. In addition, draw a trajectory passing through point $(v_0, 0)$ and another one starting at point $(0, \omega_0 r)$ Label the graph axes and indicate the direction of movement of the system for each drawn trajectory}{0.60}

Let's consider the power that is consumed during the movement of the ring.

\problembox{B4}{Calculate the instantaneous power that is consumed when there is only angular velocity $\omega$ $(v = 0)$, and separately when only linear velocity $v$ $(\omega = 0)$ is present.}{0.10}

\problembox{B5}{For given $v$ and $\omega$, calculate the instantaneous power $P$, which is spent on friction at a given time. Give your answer as an integral with a dimensionless variable.}{0.60}

\problembox{B6}{Let us assume that the ring was given a certain initial kinetic energy $E_0$. What should be the ratio $a_0 = \dfrac{v_0}{\omega_0 r}$ at which the ring will move for the maximum time? Hint: Try to answer the previous point using only $E_0$ and $a_0$ (and other data from this point), eliminating $v$ and $\omega$ from the equation}{1.20}

\problembox{B7}{What is the maximum time of motion at the initial energy $E_0$?}{0.50}

\section*{Heterogeneous system}

The system of equations that we received at the beginning of the problem is called homogeneous because the system starting to move from the point $(0,0)$ will remain at rest, that is, there is no “current” term in the equations that creates movement. Let's introduce such a term into one of the equations by considering a similar physical problem: Now we put the same ring on an inclined plane with an angle $\alpha$ with the same friction coefficient $\mu$.

\problembox{C1}{Rewrite the equations of motion from point A3 so that they fit the new condition.}{0.60}

\problembox{C2}{Given the initial $\omega_0$ and $v_0 = 0$, draw all possible families of trajectories of the ring in coordinates $(v, \omega r)$ (for each type of curve, draw your own graph).  Specify the following components: a) relevant parameter values; b) end points (to which the trajectories arrive in finite or infinite time) in the plane $(v, \omega r)$. Here it is enough to write for each component that it tends to zero / tends to infinity / is equal to or tends to some positive value. Label the axes of the graph and indicate the directions of movement of the system for each drawn trajectory}{2.00}

\vspace{1em}
\noindent\textit{Source:} \url{https://pho.rs/p/237}

\end{document}